\chapter{Технологическая часть}

В технологической части будет обоснован выбор программных средств для реализации разработанного комбинированного метода решения задачи коммивояжёра, описаны выбранные форматы входных и выходных данных. Помимо этого будет реализован сам метод, ПО для его демонстрации и взаимодействия с пользователем, проведено модульное тестирование, продемонстрирована работоспособность метода.

\section{Обоснования выбора программных средств реализации}

\subsection*{Формат входных и выходных данных}

Входными данными метода является взвешенный граф, описывающий конфигурацию задачи коммивояжёра. В качестве вводного формата для графа был выбран формат TSPLIB~\cite{tsplib}, так как:

\begin{itemize}
\item TSPLIB является стандартом для тестирования алгоритмов решения задачи коммивояжёра и имеет множество готовых конфигураций, по которым, в том числе, можно сравнить метод с другими;
\item формат поддерживает различные типы весов рёбер и графов;
\item предоставляет обширную библиотеку готовых тестовых задач различной размерности;
\end{itemize}
Выходными данными являются:
\begin{itemize}
	\item \textbf{гамильтонов цикл} -- последовательность вершин в порядке обхода, начинающаяся и заканчивающаяся в одной и той же вершине. Форматом вывода для цикла может служить строковая последовательность меток узлов графа или графическая визуализация;
	\item \textbf{длина найденного маршрута} -- сумма весов всех рёбер, входящих в найденный цикл. В зависимости от типа графа, может быть либо целым, либо вещественным числом.
\end{itemize}

\subsection*{Язык программирования}

В качестве основного языка для реализации разрабатываемого метода был выбран язык Go~\cite{go}, так как:

\begin{itemize}
	\item средств языка достаточно для реализации всех спроектированных алгоритмов и структур данных;
	\item в стандартной библиотеке присутствуют необходимые структуры данных (массивы, хеш-таблицы) и алгоритмы (сортировка, генерация случайных чисел);
	\item язык позволяет разделять программы на модули и описывать интерфейсы между ними.
\end{itemize}

\section{Особенности реализации}

\subsection*{Муравей}

В листинге~\ref{lst:ant} представлена структура муравья, хранящая индивидуальные параметры $\alpha$ и $\beta$, что обеспечивает гетерогенность колонии. Табу-список реализован с помощью встроенной хеш-таблицы go, а маршрут -- срезом.

\refstepcounter{myListingCounter}
\begin{longtable}{|p{\linewidth}|}
	\multicolumn{1}{l}{\textbf{Листинг \thechapter.\arabic{myListingCounter} -- Структура муравья}} \\
	\hline
	\endfirsthead
	\hline
	\multicolumn{1}{|c|}{\textbf{Продолжение листинга \thechapter.\arabic{myListingCounter}}} \\
	\hline
	\endhead
	\hline
	\multicolumn{1}{|r|}{} \\[-4pt]
	\hline
	\endfoot
	\hline
	\endlastfoot
	\begin{minipage}{\linewidth}
		\begin{lstlisting}[language=go, numbers=left]
type HeteroAnt struct {
	// Configuration parameters
	alpha               float64
	beta                float64
	...
	// Runtime state (initialized in Prepare)
	...
	current *graph.Vertex
	path    []*graph.Vertex
	visited map[*graph.Vertex]struct{}
	...
}
		\end{lstlisting}
	\end{minipage}
\end{longtable}
\label{lst:ant}

Алгоритм построения маршрута представлен в листинге~\ref{lst:path} : на каждом шаге муравей выбирает следующую вершину согласно вероятностному правилу, учитывающему его индивидуальные $\alpha$ и $\beta$.

\refstepcounter{myListingCounter}
\begin{longtable}{|p{\linewidth}|}
	\multicolumn{1}{l}{\textbf{Листинг \thechapter.\arabic{myListingCounter} -- Цикл поиска пути муравьём}} \\
	\hline
	\endfirsthead
	\hline
	\multicolumn{1}{|c|}{\textbf{Продолжение листинга \thechapter.\arabic{myListingCounter}}} \\
	\hline
	\endhead
	\hline
	\multicolumn{1}{|r|}{} \\[-4pt]
	\hline
	\endfoot
	\hline
	\endlastfoot
	\begin{minipage}{\linewidth}
		\begin{lstlisting}[language=go, numbers=left]
func (a *HeteroAnt) Run() error {
	...
	for !a.step() {
		// Continue until step returns true (no more vertices)
	}
	a.done = true
	return nil
}

func (a *HeteroAnt) step() bool {
	next := a.chooseNext.ChooseNext(a)
	if next == nil {
		return true // Path complete
	}
	... // обновление пути и текущей вершины
	return false
}

func (s *PathClassicStrategy) ChooseNext(ant ant.AntView) *graph.Vertex {
	...
	g.ForEachSource(current, func(edge *graph.Edge) bool {
		t := edge.Target()
		if ant.Visited(t) {
			return false
		}
		w := math.Pow(1./edge.Weight(), ant.Beta())
		p := math.Pow(pm.Get(edge), ant.Alpha())
		probability := w * p
		accumulatedProbability += probability
		candidates = append(candidates, NextCandidate{
			v:           t,
			probability: probability,
		})
		return false
	})
	...
	r := rand.Float64() * accumulatedProbability
	
	cumulative := 0.
	for _, candidate := range candidates {
		cumulative += candidate.probability
		if r <= cumulative {
			return candidate.v
		}
	}
	...
}
		\end{lstlisting}
	\end{minipage}
\end{longtable}
\label{lst:path}

После завершения построения всех маршрутов, муравьи распыляют феромон на пройденные рёбра согласно алгоритму, представленному на листинге~\ref{lst:phero}. Вызов метода происходит из управляющей структуры -- колонии муравьев.

\refstepcounter{myListingCounter}
\begin{longtable}{|p{\linewidth}|}
	\multicolumn{1}{l}{\textbf{Листинг \thechapter.\arabic{myListingCounter} -- Распыление феромона}} \\
	\hline
	\endfirsthead
	\hline
	\multicolumn{1}{|c|}{\textbf{Продолжение листинга \thechapter.\arabic{myListingCounter}}} \\
	\hline
	\endhead
	\hline
	\multicolumn{1}{|r|}{} \\[-4pt]
	\hline
	\endfoot
	\hline
	\endlastfoot
	\begin{minipage}{\linewidth}
		\begin{lstlisting}[language=go, numbers=left]
func (*ApplyClassicStrategy) ApplyPheromone(ant ant.AntView) {
	...
	delta := ant.PheromoneMultiplier() / ant.Score()
	for i := 0; i < len(path)-1; i++ {
		s, t := path[i], path[i+1]
		e, _ := g.Edge(s, t)
		pm.Add(e, delta)
	}
	wrapE, _ := g.Edge(path[len(path)-1], path[0])
	pm.Add(wrapE, delta)
}
		\end{lstlisting}
	\end{minipage}
\end{longtable}
\label{lst:phero}

\subsection*{Локальная оптимизация}

Для повышения качества найденных решений и ускорения сходимости в методе может применяться локальная оптимизация маршрутов. На листинге~\ref{lst:2-opt} представлена реализация алгоритма 2-opt, которая применяется к отобранным родителям перед формированием нового поколения.

\refstepcounter{myListingCounter}
\begin{longtable}{|p{\linewidth}|}
	\multicolumn{1}{l}{\textbf{Листинг \thechapter.\arabic{myListingCounter} -- 2-opt оптимизация}} \\
	\hline
	\endfirsthead
	\hline
	\multicolumn{1}{|c|}{\textbf{Продолжение листинга \thechapter.\arabic{myListingCounter}}} \\
	\hline
	\endhead
	\hline
	\multicolumn{1}{|r|}{} \\[-4pt]
	\hline
	\endfoot
	\hline
	\endlastfoot
	\begin{minipage}{\linewidth}
		\begin{lstlisting}[language=go, numbers=left]
func (s TwoOptLocalOptimisation) Optimise(path []graph.Vertex, g *graph.Graph) {
	n := len(path)
	improved := true
	for improved {
		improved = false
		for i := 0; i < n-2; i++ {
			for j := i + 2; j < n-1; j++ {
				e1, _ := g.Edge(path[i], path[i+1])
				e2, _ := g.Edge(path[j], path[j+1])
				
				if s.shouldSwap(e1, e2, g) {
					s.swapEdges(path, i, j)
					improved = true
					break
				}
			}
			if improved {
				break
			}
		}
	}
}

func (s *TwoOptLocalOptimisation) shouldSwap(e1, e2 *graph.Edge, g *graph.Graph) bool {
	v1, v2 := e1.Source(), e1.Target()
	u1, u2 := e2.Source(), e2.Target()
	// Текущая длина двух рёбер
	curWeight := e1.Weight() + e2.Weight()
	// Новая длина после перестановки
	newWeight := g.Edge(v1, u1).Weight() + g.Edge(v2, u2).Weight()
	return newWeight < curWeight // замена сокращает маршрут
}

func (s TwoOptLocalOptimisation) swapEdges(path []graph.Vertex, i, j int) {
	// Переворот участка между i+1 и j
	for start, end := i+1, j; start < end; start, end = start+1, end-1 {
		path[start], path[end] = path[end], path[start]
	}
}
		\end{lstlisting}
	\end{minipage}
\end{longtable}
\label{lst:2-opt}

\subsection*{Колония}

Колония (листинг~\ref{lst:colony}) является центральным управляющим компонентом метода. Она хранит популяцию муравьёв, карту феромонов, а также настройки и стратегии генетического алгоритма. Поля \texttt{parentSelect}, \texttt{crossover}, \texttt{mutation} представляют собой стратегии селекции, скрещивания и мутации соответственно. Срез \texttt{observers} содержит зарегистрированных наблюдателей для сбора статистики.

\refstepcounter{myListingCounter}
\begin{longtable}{|p{\linewidth}|}
	\multicolumn{1}{l}{\textbf{Листинг \thechapter.\arabic{myListingCounter} -- Структура колонии}} \\
	\hline
	\endfirsthead
	\hline
	\multicolumn{1}{|c|}{\textbf{Продолжение листинга \thechapter.\arabic{myListingCounter}}} \\
	\hline
	\endhead
	\hline
	\multicolumn{1}{|r|}{} \\[-4pt]
	\hline
	\endfoot
	\hline
	\endlastfoot
	\begin{minipage}{\linewidth}
		\begin{lstlisting}[language=go, numbers=left]
type HeteroAntColony struct {
	// Стратегии
	parentSelect      ParentSelectionStrategy
	crossover         CrossoverStrategy
	mutation          MutationStrategy
	// Параметры муравьиного алгоритма
	defaultAlpha        float64
	defaultBeta         float64
	pheromoneMultiplier float64
	evaporationRate     float64
	initialPheromone    float64
	// Параметры генетического алгоритма
	generationCount  uint
	colonySize       uint
	generationPeriod uint
	parentCount      uint
	...
	best  *ant.HeteroAnt
	score float64
	...
	observers []ColonyObserver // Наблюдатели
}
		\end{lstlisting}
	\end{minipage}
\end{longtable}
\label{lst:colony}

Основной цикл метода, представленный в листинге~\ref{lst:method}, реализует итерационный процесс:

\begin{itemize}
	\item на каждой итерации все муравьи строят маршруты;
	\item каждые $w$ итераций (период адаптации \texttt{generationPeriod}) происходит отбор родителей заданной стратегией, к которым применяется локальная оптимизация 2-opt;
	\item после построения маршрутов выполняется испарение феромонов;
	\item затем муравьи распыляют феромон на пройденные рёбра;
	\item наблюдатели уведомляются о текущем состоянии для сбора статистики;
	\item если на данной итерации выполнялась селекция, формируется новое поколение, где последовательно применяются скрещивание и мутация, переданные при конфигурации.
\end{itemize}

\refstepcounter{myListingCounter}
\begin{longtable}{|p{\linewidth}|}
	\multicolumn{1}{l}{\textbf{Листинг \thechapter.\arabic{myListingCounter} -- Основной цикл метода}} \\
	\hline
	\endfirsthead
	\hline
	\multicolumn{1}{|c|}{\textbf{Продолжение листинга \thechapter.\arabic{myListingCounter}}} \\
	\hline
	\endhead
	\hline
	\multicolumn{1}{|r|}{} \\[-4pt]
	\hline
	\endfoot
	\hline
	\endlastfoot
	\begin{minipage}{\linewidth}
		\begin{lstlisting}[language=go, numbers=left]
func (c *HeteroAntColony) Run() error {
	for i := uint(0); i < c.colonySize; i++ {...} // Инициализация первого поколения
	// Основной цикл
	for gen := uint(0); gen < c.generationCount; gen++ {
		for _, a := range c.ants {...a.Run()...} // Построение путей муравьями
		var parents []ant.AntView
		if gen%c.generationPeriod == 0 { // Отбор
			...
			parents = c.parentSelect.SelectParents(views, c.parentCount)
			for _, p := range parents {...a.LocalOptimisation(c.localOptimisation)...} // Локальная оптимизация
			...
		}
		// Поиск более лучшего пути среди найденных
		for _, a := range c.ants {
			if c.best == nil || bestInGen.Score() < c.best.Score() {
				c.best = a
			}
		}
		c.evaporatePheromones() // Испарение 
		for _, a := range c.ants {...a.ApplyPheromone()...} // Нанесение феромонов 
		dto := &ColonyObserverDTO{...}
		for _, obs := range c.observers { // Уведомление наблюдателей
			obs.Observe(dto)
		}
		if parents != nil {
			c.prepareNextGeneration(parents) // Обновление поколения генетическим алгоритмом
		}
	}
	return nil
}
		\end{lstlisting}
	\end{minipage}
\end{longtable}
\label{lst:method}

Вспомогательные функции метода, представленны в листинге~\ref{lst:method-side}.

\refstepcounter{myListingCounter}
\begin{longtable}{|p{\linewidth}|}
	\multicolumn{1}{l}{\textbf{Листинг \thechapter.\arabic{myListingCounter} -- Основной цикл метода}} \\
	\hline
	\endfirsthead
	\hline
	\multicolumn{1}{|c|}{\textbf{Продолжение листинга \thechapter.\arabic{myListingCounter}}} \\
	\hline
	\endhead
	\hline
	\multicolumn{1}{|r|}{} \\[-4pt]
	\hline
	\endfoot
	\hline
	\endlastfoot
	\begin{minipage}{\linewidth}
		\begin{lstlisting}[language=go, numbers=left]
func (c *HeteroAntColony) evaporatePheromones() {
	c.g.ForEachEdge(func(e *graph.Edge) bool {
		c.pm.Update(e, c.pm.Get(e)*(1 - c.evaporationRate))
	})
}
func (c *HeteroAntColony) prepareNextGeneration(parents []ant.AntView) {
	offspring := make([]*ant.HeteroAnt, 0, c.colonySize) // Скрещивание
	for len(offspring) < int(c.colonySize) {
		p1, p2 := parents[rand.Intn(len(parents) - 1)], parents[rand.Intn(len(parents) - 1)]
		offspring = append(offspring, c.crossover.Crossover(p1, p2)...)
	}
	for i := 0; uint(i) < c.colonySize; i++ { // Мутация
		offspring[i] = c.mutation.Mutate(offspring[i])
	}
}
		\end{lstlisting}
\end{minipage}
\end{longtable}
\label{lst:method-side}

\subsection*{Стратегии}

В листинге~\ref{lst:strategies} представлены интерфейсы стратегий генетического алгоритма, которые используются колонией для эволюции популяции. Как видно, параметры  мутации хранятся напрямую в её структуре и передаются при создании её экземпляра, что даёт возможность не включать всевозможные параметры стратегий в основную структуру колонии.

\refstepcounter{myListingCounter}
\begin{longtable}{|p{\linewidth}|}
	\multicolumn{1}{l}{\textbf{Листинг \thechapter.\arabic{myListingCounter} -- Интерфейсы стратегий}} \\
	\hline
	\endfirsthead
	\hline
	\multicolumn{1}{|c|}{\textbf{Продолжение листинга \thechapter.\arabic{myListingCounter}}} \\
	\hline
	\endhead
	\hline
	\multicolumn{1}{|r|}{} \\[-4pt]
	\hline
	\endfoot
	\hline
	\endlastfoot
	\begin{minipage}{\linewidth}
		\begin{lstlisting}[language=go, numbers=left]
type ParentSelectionStrategy interface {
	SelectParents(ants []ant.AntView, n uint) []ant.AntView
}

type CrossoverStrategy interface {
	Crossover(p1, p2 ant.AntView) []*ant.HeteroAnt
}

type MutationStrategy interface {
	Mutate(ant ant.AntView) *ant.HeteroAnt
}
		\end{lstlisting}
	\end{minipage}
\end{longtable}
\label{lst:strategies}

Как пример реализации одного из интерфейсов, в листинге~\ref{lst:gauss} приведена реализация одного из видов мутации -- гауссовской.

\refstepcounter{myListingCounter}
\begin{longtable}{|p{\linewidth}|}
	\multicolumn{1}{l}{\textbf{Листинг \thechapter.\arabic{myListingCounter} -- Гауссова мутация}} \\
	\hline
	\endfirsthead
	\hline
	\multicolumn{1}{|c|}{\textbf{Продолжение листинга \thechapter.\arabic{myListingCounter}}} \\
	\hline
	\endhead
	\hline
	\multicolumn{1}{|r|}{} \\[-4pt]
	\hline
	\endfoot
	\hline
	\endlastfoot
	\begin{minipage}{\linewidth}
		\begin{lstlisting}[language=go, numbers=left]
type GaussMutationStrategy struct {
	sigma float64
	mu    float64
}

var _ colony.MutationStrategy = (*GaussMutationStrategy)(nil)

func normalRand(sigma, mu float64) float64 {
	return rand.NormFloat64()*sigma + mu
}

func NewGaussMutationStrategy(sigma, mu float64) *GaussMutationStrategy {
	return &GaussMutationStrategy{
		sigma: sigma,
		mu:    mu,
	}
}

func (s *GaussMutationStrategy) Mutate(a ant.AntView) *ant.HeteroAnt {
	alpha := normalRand(s.sigma, s.mu)
	beta := normalRand(s.sigma, s.mu)
	
	return ant.NewHeteroAnt(
	a.Alpha()+alpha,
	a.Beta()+beta,
	a.PheromoneMultiplier(),
	a.PathStrategy(),
	a.PheromoneApplyStrategy(),
	)
}
		\end{lstlisting}
	\end{minipage}
\end{longtable}
\label{lst:gauss}


\subsection*{Наблюдатели}

Для сбора статистики и мониторинга работы алгоритма в методе реализован паттерн <<Наблюдатель>>. Интерфейс \texttt{ColonyObserver} (листинг~\ref{lst:observer}) содержит метод \texttt{Observe}, принимающий структуру данных \texttt{ColonyObserverDTO}. Она инкапсулирует текущее состояние колонии: популяцию муравьёв, карту феромонов, лучшее найденное решение и номер текущей итерации и позволяет наблюдателям агрегировать эти данные для статистики.

\refstepcounter{myListingCounter}
\begin{longtable}{|p{\linewidth}|}
	\multicolumn{1}{l}{\textbf{Листинг \thechapter.\arabic{myListingCounter} -- Интерфейс наблюдателя}} \\
	\hline
	\endfirsthead
	\hline
	\multicolumn{1}{|c|}{\textbf{Продолжение листинга \thechapter.\arabic{myListingCounter}}} \\
	\hline
	\endhead
	\hline
	\multicolumn{1}{|r|}{} \\[-4pt]
	\hline
	\endfoot
	\hline
	\endlastfoot
	\begin{minipage}{\linewidth}
		\begin{lstlisting}[language=go, numbers=left]
type ColonyObserverDTO struct {
	C     *HeteroAntColony
	Ants  []ant.AntView
	Pm    *pheromone.PheromoneMap
	G     *graph.Graph
	Best  *ant.HeteroAnt
	Score float64
	
	Generation uint
}

type ColonyObserver interface {
	Observe(dto *ColonyObserverDTO)
}
		\end{lstlisting}
	\end{minipage}
\end{longtable}
\label{lst:observer}

В листинге~\ref{lst:runtime} представлен пример наблюдателя -- \texttt{TimeObserver}, предназначенный для замера времени выполнения. Он сохраняет временные метки для каждой итерации, что позволяет впоследствии проанализировать производительность алгоритма.

\refstepcounter{myListingCounter}
\begin{longtable}{|p{\linewidth}|}
	\multicolumn{1}{l}{\textbf{Листинг \thechapter.\arabic{myListingCounter} -- Наблюдатель времени выполнения}} \\
	\hline
	\endfirsthead
	\hline
	\multicolumn{1}{|c|}{\textbf{Продолжение листинга \thechapter.\arabic{myListingCounter}}} \\
	\hline
	\endhead
	\hline
	\multicolumn{1}{|r|}{} \\[-4pt]
	\hline
	\endfoot
	\hline
	\endlastfoot
	\begin{minipage}{\linewidth}
		\begin{lstlisting}[language=go, numbers=left]
type TimeObserver struct {
	startTime time.Time
	endTime   time.Time
	runs      []RunTime
}
...
func (o *TimeObserver) Observe(dto *colony.ColonyObserverDTO) {
	if len(o.runs) == 0 {
		o.runs = append(o.runs, RunTime{
			Run:    dto.Generation,
			Moment: time.Now(),
			Time:   time.Since(o.startTime),
		})
		return
	}
	lastRun := o.runs[len(o.runs)-1]
	o.runs = append(o.runs, RunTime{
		Run:    dto.Generation,
		Moment: time.Now(),
		Time:   time.Since(lastRun.Moment),
	})
}
		\end{lstlisting}
	\end{minipage}
\end{longtable}
\label{lst:runtime}

\section{Тестирование}

Модульное и интеграционное тестирование приложения проводилось с помощью стандартного пакета Go -- testing, и стороннего пакета testify.

В качестве меры качества покрытия кода использовалась метрика стандартной утилиты cover входящей в состав Go, которая измеряет покрытие кода как процент от выполненных хотя бы единожды выражений~\cite{go-cover}.

Тестами были покрыты модули стратегий, структур данных, а также части ядра, посвященные вычислениям. В таблице~\ref{tbl:tests} представлено покрытие кода тестами по модулям.

\begin{longtable}{|p{0.4\textwidth}|p{0.3\textwidth}|}
	\caption{Результаты модульного тестирования}\label{tbl:tests} \\
	\hline
	\textbf{Модуль} & \textbf{Процент покрытия тестами} \\
	\hline
	\endfirsthead
	\caption*{Продолжение таблицы~\ref{tbl:tests}} \\
	\hline
	\textbf{Модуль} & \textbf{Процент покрытия тестами} \\
	\hline
	\endhead
	\hline
	\endfoot
	\endlastfoot
	Ядро & 45.6\% \\
	\hline
	Структуры данных & 72.9\% \\
	\hline
	Стратегии & 82.3\% \\
	\hline
	Наблюдатели & 0\% \\
	\hline
	\textbf{Всего} & \textbf{44.3\%} \\
	\hline
\end{longtable}

\section{Пример работы ПО}

Для демонстрации работы разработанного метода был разработан web-интерфейс, позволяющий вводить параметры и визуализировать полученное решение, характеристики алгоритма.

В качестве примера работы программы взято решение задачи с входным графом \texttt{ulysses22.tsp}. На рисунке~\ref{fig:params} представлен экран ввода графа и параметров метода.

\begin{figure}[H]
	\centering
	\includegraphics[width=0.6\textwidth]{examples/params}
	\caption{Экран ввода параметров}
	\label{fig:params}
\end{figure}

На рисунке~\ref{fig:stats} представлен экран с выводом информации о характеристиках входного графа, а также характеристиках решения и статистики запуска. Решение задачи в виде гамильтонова цикла графа изображено на рисунке~\ref{fig:graphs}. Там же отображена карта феромонов графа.

\begin{figure}[H]
	\centering
	\includegraphics[width=0.6\textwidth]{examples/stats}
	\caption{Экран вывода информации о графе и решении}
	\label{fig:stats}
\end{figure}

\begin{figure}[H]
	\centering
	\includegraphics[width=0.6\textwidth]{examples/graphs}
	\caption{Экран визуализации решения задачи и карты феромонов}
	\label{fig:graphs}
\end{figure}

На рисунке~\ref{fig:evolve} представлен экран визуализации изменения средних значений коэффициентов $\alpha$ и $\beta$, а также эволюция длин лучшего пути по поколениям муравьёв.

\begin{figure}[H]
	\centering
	\includegraphics[width=0.6\textwidth]{examples/evolve}
	\caption{Экран визуализации изменения параметров и лучшего пути}
	\label{fig:evolve}
\end{figure}

\section*{Вывод}

В результате был обоснован выбор языка программирования Go как средства реализации предложенного метода. Приведены ключевые фрагменты реализации, демонстрирующие: структуру муравья с гетерогенными параметрами $\alpha$ и $\beta$, вероятностное правило выбора следующей вершины, локальную оптимизацию 2-opt, основной цикл колонии с периодической эволюцией популяции, интерфейсы стратегий генетического алгоритма и реализацию паттерна «Наблюдатель» для сбора статистики. Проведено тестирование программной реализации, покрытие кода составило 42,2\%. Представлен пример работы на тестовой задаче \texttt{ulysses22.tsp} с визуализацией решения и эволюции параметров.


\clearpage
